% Section: Introduction

\section{Introduction}
\label{sec:intro}

Attribution graphs provide a structured view of which learned features mediate a model's prediction for a given target logit~\cite{ameisen2025circuittracing}. In practice, interpreting these graphs remains bottlenecked by extensive manual analysis: identifying influential features, inferring their semantics from corpus activations, and assembling them into a legible subgraph. A typical workflow requires researchers to manually inspect activation patterns across dozens of context variations to determine whether a given feature encodes semantic content (e.g., `France'), relational structure (e.g., `capital-of'), or syntactic patterns (e.g., `says X'). As noted in the Neuronpedia podcast, analyzing a single prompt takes ``approximately 2 hours for an experienced circuit tracer.'' This process demands deep expertise and produces inconsistent taxonomies across researchers. As interpretability scales to larger model families and more complex behaviors, manual circuit analysis becomes increasingly infeasible.

We make three empirical claims about automated circuit interpretation. First, concept-aligned grouping based on behavioral signatures exhibits higher functional coherence than geometric clustering (C1): our supernodes show 2.3$\times$ higher peak-token consistency (0.425 vs 0.183) and 5.8$\times$ higher activation-pattern similarity (0.762 vs 0.130) than cosine-based grouping (Table~\ref{tab:baselines}). Second, behavioral grouping enables interpretability-oriented compression (C2): subgraphs preserve 83\% of explanatory coverage (Completeness) while accepting 54\% end-to-end influence routing (Replacement) as a deliberate trade-off for legibility (Table~\ref{tab:graph-subgraph}). Third, entity-swap tests reveal a layerwise computational hierarchy (C3): early-layer features transfer robustly across entity substitutions (64\% transfer, mean layer 6.3), while late-layer Say-X features specialize for output promotion (mean layer 16.4), supporting a backbone-and-specialization model of transformer computation (Table~\ref{tab:cross-prompt}).

We evaluate these claims across four prompt families: geographic `capitals' circuits with entity swaps (Dallas$\rightarrow$Austin/Texas, Oakland$\rightarrow$Sacramento/California), a sports-name circuit (Michael Jordan$\rightarrow$Basketball), an antonymy pair (small$\rightarrow$opposite), and an anatomy pair (muscle$\rightarrow$diaphragm). Our evidence combines graph-level metrics (Neuronpedia's Replacement and Completeness), behavioral coherence metrics (peak-token consistency, activation-pattern similarity, sparsity consistency), and transfer analysis via entity substitution. Results are descriptive rather than inferential; we document evaluation procedures and flag exploratory elements to support transparent calibration (see table notes for full details).

Our work builds on recent mechanistic-interpretability infrastructure: attribution graphs formalize feature→logit pathways, and platforms such as Neuronpedia provide feature cards and graph tooling at scale~\cite{ameisen2025circuittracing,neuronpedia2025}. Sparse Autoencoders (SAE) and Cross-Layer Transcoders (CLT) expose sparse, often monosemantic features that make replacement-model analyses tractable~\cite{templeton2024scaling,bricken2023monosemanticity,ameisen2025circuittracing}. Building on our earlier Prompt Rover framing of ``semantic navigation''~\cite{birardi2025insight} and its accompanying tool~\cite{birardi2025promptrover}, we operationalize automated concept-hypothesis generation and cross-prompt behavioral measurement, with transparent rule-based grouping and naming. The novelty is not a new distance or clustering algorithm; it is the shift in the unit of interpretation---from geometric similarity to functional behavior measured across systematically varied contexts. Metaphorically: project multiple lights onto a complex object (probe prompts); each light casts a different shadow (activation pattern); comparing these shadows reveals the object's underlying conceptual structure.

Following best practices for scientific integrity, we frame probe prompting as an interpretability method that prioritizes legibility and behavioral coherence. Our quantitative results are descriptive rather than inferential, and our validation is correlational. Planned causal tests (ablations, steering, cross-circuit transfer) are outlined in Section~\ref{sec:discussion} as targets for future work. To facilitate adoption, we provide an interactive demo at \url{https://huggingface.co/spaces/Peppinob/attribution-graph-probing} where researchers can upload attribution graphs and explore probe-prompted subgraphs without local installation, alongside our public codebase at \url{https://github.com/peppinob-ol/attribution-graph-probing}. Epistemic status: medium confidence in current pipeline on Gemma-2-2B/CLT-HP with short factual prompts; low confidence in broad generalization to long chain-of-thought reasoning, adversarial prompts, or larger models without adaptation.


